\section{Conclusion}
Transformers use a Softmax-based mechanism known as attention, which has a \textit{quadratic} scaling with number of tokens taken in context. This problem is particularly important since the current applications are moving towards higher number of tokens, and that there is significant interest in deploying large language models on devices with limited computational power such as smart phones. Consequently, the search for an efficient attention mechanism has become an important and active area of research. One of the promising approaches is Linear Attention (LA) with a linear time complexity of $O(ND^2)$, where $N$ is the number of tokens and $D$ the dimension per attention head. However, due to the less established nature of LA, its implementation is far less optimized compared to regular attention and there is much room for improvement. We present an extensively optimized LA implementation.

\par In summary, we propose a novel approach for calculating LA, use it to derive the analytical gradient of Attention heads, demonstrate how forward and backward passes can be calculated in $O(N)$ time, and provide the details of our optimized implementation. Our method and its implementation work hand-in-hand in the sense that the computational pattern we devise allows for a highly parallelized implementation with minimal data movement. We evaluate our implementation in an isolated setting as a standalone attention layer, and in an end-to-end setting where we train an LLM. We achieve $3.3\times$ speedup and a $3.6\times$ reduction in memory consumption compared to the state-of-the-art LA implementation~\cite{yang2023gated}. Our implementation is in format of a PyTorch library, and can be used as a plug-and-play module. By improving the efficiency of LA implementation, we are paving the path for its practical usage, which will be particularly beneficial for deploying LLMs on computationally constrained devices. 

% Transformers use a Softmax-based mechanism known as Attention, which has a quadratic scaling with number of tokens taken in context. Many approaches have proposed subquadratic mechanisms for both training and inference. This problem is particularly important since the current applications are moving towards higher number of tokens, and training can only be done in large datacenter scale systems. The search for a linearly scaling Attention mechanism has become an important and active area of research. One of the promising approaches is Linear Attention (Kernel Separation). However, current implementations had very  high memory consumption of $O(ND^{2})$, where $N$ is the number of tokens and $D$ the model dimension per head. This memory consumption makes kernel separation impractical since even small transformers would require a high amount of memory per batch. We present a solution to reduce the memory scaling of Linear Attention to $O(ND)$. In summary, we derive the analytical format of the Attention layer gradient and calculate it in linear time. By implementing this approach instead of relying on a differentiable programming library, we eliminate the need to store the variables for the backward pass, and bring down the memory consumption to $O(ND)$. We provide the details of how to implement the forward and backward pass using Factorization, and confirm the time and memory scaling of our implementation in context of an Attention layer and training an LLM. Furthermore, since Linear Attention does not directly compute the Attention matrix, dropout cannot be applied. Therefore, we have also introduced an alternative mechanism to dropout, to mimic its effect, and confirmed our mechanism's effectiveness through an ablation study.